
\subsection{What the observed state explains}
At exactly 512 actual prompt tokens, mean completed-step time rises from
13.74 to 17.81 ms when the decoder-count bin changes from 1--7 to 16--23;
the reverse repetition gives 14.07 to 17.93 ms. At 2,048 actual prompt tokens,
the 8--15 versus 16--23 bins give 24.92 versus 27.38 ms, reproduced as
24.93 versus 27.37 ms. These comparisons hold prompt work fixed, not total tokens.
For a stricter fixed-total comparison, two fixed-1,024 runs each contain 19 steps
with exactly 10 decoders. Splitting them by context sum gives mean costs
18.268 versus 18.730 ms and 18.281 versus 18.741 ms. This is a small repeatable
association at 1,034 total scheduled tokens. A five-step bin with 18 decoders
has the opposite direction, so we do not infer a monotonic causal law or a
validated independent prediction signal.

\begin{figure}[htbp]\centering
\begin{tikzpicture}
\begin{axis}[width=.82\textwidth,height=5.3cm,xlabel={Active decoders},
 ylabel={Completed step envelope (ms)},grid=major,
 legend style={at={(.03,.97)},anchor=north west},legend cell align=left]
\addplot[only marks,mark=*,mark size=1.4pt,blue] table[x=decode,y=ms,col sep=comma]{step_cost_512.csv};
\addlegendentry{Exactly 512 prefill tokens}
\addplot[only marks,mark=triangle*,mark size=1.8pt,red] table[x=decode,y=ms,col sep=comma]{step_cost_2048.csv};
\addlegendentry{Exactly 2,048 prefill tokens}
\end{axis}
\end{tikzpicture}
\caption{Actual mixed steps from the first measured fixed pair. Partial tail
chunks are excluded only from this equal-prompt-work diagnostic, not from service metrics.}
\end{figure}

The fixed-1,024 and feedback development runs agree on the complete output
sequence for 24 of 36 requests; each policy agrees on all 36 sequences between
its own repetitions. Every run still generates exactly 8,448 tokens. This is
consistent with shape-dependent numerical execution changing greedy trajectories;
it does not establish either a precision change or a task-quality difference.
